\color{unidarkgreen}{\subsection{Блок отстройки от повреждений за трансформаторами ответвлений}}
\color{black}

\begin{enumerate}[label=\arabic{section}.\arabic{subsection}.\arabic{enumi}, labelsep=4pt, leftmargin=0pt, itemindent=57pt, itemsep=0pt, parsep=5pt] % если что можно крутить itemsep

%%%%%%%%%%%%%%%%%%%%%%%%%%%%%%%%%%%%%%%%%%%%%%%%%% ОБШИЕ СВЕДЕНИЯ %%%%%%%%%%%%%%%%%%%%%%%%%%%%%%%%%%%%%%%%%%%%%%%%%%
\item Общие сведения

\begin{enumerate}[label=\arabic{section}.\arabic{subsection}.\arabic{enumi}.\arabic*, labelsep=4pt, leftmargin=0em, itemindent=65pt, parsep=0pt]

\item
Для предотвращения ложного срабатывания основной защиты при КЗ за трансформатором ответвления предусмотрен функциональный блок отстройки от повреждений за трансформаторами ответвлений (ОПТО).
В качестве измерительных органов ОПТО применяются три междуфазных прямонаправленных реле сопротивления и РНМНП-Р, блокируемое при выявлении броска намагничивающего тока.
При введенном программном ключе контроля защиты от ОПТО подготовка её цепи отключения разрешается только при срабатывании одного из реле сопротивления или РНМНП-Р, что свидетельствует о наличии КЗ на линии, а не за трансформатором ответвления.
Характеристика срабатывания РС задается уставками \textbf{«Rср РС-ОПТО»}, \textbf{«Хср РС-ОПТО»} и \textbf{«Ф1 РС-ОПТО»}.
Углы наклона характеристики срабатывания органа направленности, обеспечивающего направленность РС, задаются уставками \textbf{«Ф2 ОН-ОПТО»} и \textbf{«Ф3 ОН-ОПТО»}.
Для РС в составе ОПТО также предусмотрена возможность выреза в характеристике срабатывания для отстройки от нагрузки.
Характеристика срабатывания измерительного органа отстройки от нагрузки определяется уставками \textbf{«Rнг ОПТО»} и \textbf{«Фнг ОПТО»}.
Работа ОПТО может быть заблокирована при выявлении неисправности в цепях напряжения при введенном программном ключе \textbf{«БНН ОПТО»}.

\medskip

\begin{figure}[H]
\centering
\includegraphics[width=0.4\textwidth,height=0.4\textheight,keepaspectratio]{\fbpath/561.1491 ОПТО/img/RS_OPTO.pdf}
\captionof{figure}{Характеристика срабатывания РС в составе ОПТО}\label{pict_opto:RS_OPTO}
\end{figure}

\medskip

\begin{figure}[H]
\centering
\includegraphics[width=1\textwidth,height=1\textheight,keepaspectratio]{\fbpath/561.1491 ОПТО/img/OPTO.pdf}
\captionof{figure}{Функциональная схема отстройки от повреждений за трансформаторами ответвлений}
\label{pict_opto:OPTO}
\end{figure}

\medskip

\item
Уставки ОПТО представлены в таблице \ref{desc_opto:tbl1}.

\small
\begin{longtable}{|>{\centering\arraybackslash}m{5.3cm}|>{\centering\arraybackslash}m{3.3cm}|>{\centering\arraybackslash}m{4.2cm}|>{\centering\arraybackslash}m{1.8cm}|>{\centering\arraybackslash}m{1cm}|}
% Первая страница: полный заголовок
\caption{Уставки ОПТО\hfill\vspace{-0.5\baselineskip}}\label{desc_opto:tbl1}\\
\hline
\rowcolor{gray!30}
Параметр & Условное обозначение на схеме & Значение/ Диапазон & Единица измерения & Шаг \\
\hline
\endfirsthead

% Промежуточные страницы: «Продолжение таблицы»
\caption*{\hspace{3pt}\emph{Продолжение таблицы \ref{desc_opto:tbl1}}\hfill\ }\vspace{-0.5\baselineskip} \\
\hline
\rowcolor{gray!30}
Параметр & Условное обозначение на схеме & Значение/ Диапазон & Единица измерения & Шаг \\
\hline
\endhead

%===m>ОПТО|-
%===m
\end{longtable}
\normalsize

\end{enumerate}

\end{enumerate}
